\documentclass{article}
\usepackage[left=2cm, right=2cm, top=2cm, bottom=2cm]{geometry}
\usepackage{graphicx}
\usepackage{amsmath}
\usepackage{amssymb}
\usepackage{amsthm}
\usepackage{fancyhdr}
\usepackage{verbatim}
\usepackage{listings}
\usepackage{xcolor}
\usepackage{pdfpages}
\usepackage{cancel}
\usepackage{physics}
\usepackage{esint}
\usepackage{braket}

\lstset{
    language=C++,
    basicstyle=\ttfamily\footnotesize,
    keywordstyle=\color{blue},
    commentstyle=\color{green},
    stringstyle=\color{red},
    numbers=left,
    numberstyle=\tiny\color{gray},
    frame=single,
    breaklines=true,
    captionpos=b,
}


\title{MTH 446: Lecture \# 10}
\author{Cliff Sun}

\newtheorem{theorem}{Theorem}[section]
\newtheorem{lemma}[theorem]{Lemma}
\newtheorem{definition}[theorem]{Definition}
\newtheorem{conjecture}[theorem]{Conjecture}
\newtheorem{proposition}[theorem]{Proposition}
\newtheorem{corollary}[theorem]{Corollary}
\newtheorem{one minute paper}[theorem]{One Minute Paper}

\pagestyle{fancy}
\lhead{\textbf{\thepage}\ \ \nouppercase{\rightmark}}
\chead{MTH 446: Lecture \# 10}
\rhead{Cliff Sun}

\begin{document}

\maketitle

\section*{Lecture Span}
\begin{enumerate}
    \item Limits of complex functions
\end{enumerate}
 
\section*{Recap}

You may have to extend functions to the extended complex plane. An example

\begin{equation}
    \tilde{f}(z) = \begin{cases}
        1/z & z \neq 0 \lor \infty \\
        0 & z = \infty \\
        \infty & z=0
    \end{cases}
\end{equation}

Use algebraic properties of limits. For the third case, prove that $1/f(z) = z \rightarrow 0$. Clearly, this is true.

Also, we need to show that $z \rightarrow z_0$, where $z_0 \in G = \text{dom}(f)$. In other words, $z_0$ needs to be in the domain of $f$. 

For the second case, $|z| > 1/\delta \implies |1/z - 0| = 1/|z| < \epsilon \iff |z| > 1/\epsilon$. Therefore, let $\delta = \epsilon$, then we see that 
\begin{equation}
    \lim_{z \rightarrow 0} (1/z) = 0
\end{equation}

Therefore, $\tilde{f} \in C(\overline{\mathbb{C}})$ (cont. in the extended complex plane). 

\subsection*{Example}

(Mobius transformation) Defined by 
\begin{equation}
    M(z) = \frac{az + b}{az + d}:  a,b,c,d \in \mathbb{C}
\end{equation}
\begin{equation}
    \det\begin{pmatrix}
        a & b \\
        c & d
    \end{pmatrix} \neq 0
\end{equation}

Therefore, if $M(z) = 1/z = (0 \cdot z + 1)/(1 \cdot z + 0)$ which is 

\begin{equation}
    \begin{pmatrix}
        0 & 1 \\
        1 & 0 
    \end{pmatrix}
\end{equation}

Extend $M$ to the extended complex plane. Consider 

\begin{equation}
    M_1 \circ M_2 = M_1(M_2(z))
\end{equation}

With $\det M_1 \neq 0$ and $\det M_2 \neq 0$. Let $c \neq 0$, then we can extend $M$ to the extended complex plane:

\begin{equation}
    \tilde{M}(z) = \begin{cases}
        \frac{az + b}{cz + d} & z \neq \infty, z \neq -d/c \\
        \infty & z = -d/c \\
        a/z & z = \infty
    \end{cases}
\end{equation}

Case 1: if $z \neq \infty$ and $z \neq -d/c$. Then $az + b$ and $cz + d$ are in $C(\overline{\mathbb{C}})$. Moreover, $cz + d \neq 0$. Then, by the algebraic properties of cont. functions, then $\tilde{M}(z)$ is cont. 

Case 2: if $z = -d/c$. Firstly, $-d/c \in \overline{\mathbb{C}}$. Next, write

\begin{equation}
    \lim_{z \rightarrow -d/c} \tilde{M}(z) = \infty \iff \lim_{z \rightarrow -d/c} 1/\tilde{M}(z) = 0
\end{equation}

Then 

\begin{align}
    \lim 1/\tilde{M} &= \lim \frac{cz + d}{az + b} \\
    \frac{cz + d}{az + b} \rightarrow \frac{0}{- ad/c + b} = 0 \iff ad - bc \neq 0
\end{align}

But, this is the definition of the Mobius transformation, therefore, $1/\tilde{M}(z) \rightarrow 0 \iff \tilde{M}(z) \rightarrow \infty$.

Case 3: $z_0 \rightarrow \infty$. Then, evaluate 

\begin{equation}
    \lim \frac{az + b}{cz + d} = \lim \frac{a + b/z}{c + d/z} \rightarrow a/c
\end{equation}

Here, we used the algebraic properties of cont. functions.  

Next, if $c = 0$, then $ad - bc = ad \neq 0$ which means that $a \neq 0$ and $d \neq 0$. Therefore, 

\begin{equation}
    M(z) = a/d \cdot z + b/d
\end{equation}

Then 

\begin{equation}
    M(z) = A z + B, A \neq 0
\end{equation}

We define this transformation otno the extended complex plane. in other words, 

\begin{equation}
    \tilde{M}(z) = \begin{cases}
        Az + B  & z \neq \infty\\
        \infty & z = \infty
    \end{cases}
\end{equation}

Case 1: obvious. 

Case 2:

\begin{align}
    \lim 1/M(z) &= \lim \frac{1}{Az + B}  = \lim \frac{1/z}{A + B/z} \rightarrow 0/A = 0
\end{align}

Therefore, $\lim \tilde{M}(z) = \infty$ and $\tilde{M}(\infty) = \infty$. Therefore, $\tilde{M} \in \overline{\mathbb{C}}$.

\textbf{Midterm:}
\begin{enumerate}
    \item 5 problems 20 points each, no cheat sheet. 
    \item problem 1: general information about complex numbers. Conjugate, modulus. $|z|^2 = z \overline{z}$. 
    \item problem 2: roots of complex number. Principal root
    \item problem 3: Problem of limits
    \item problem 4: Problem related to Demoire formula
    \item problem 5: problem related to open set, closed set, interior, boundary (D intersects with outside and inside for all $\epsilon$). Find closure. Accumulation points. 
\end{enumerate}

\section*{Complex Derivative}

\begin{definition}
    $f$ is complex differentiable at $z = z_0$ if the following exists:
    \begin{equation}
        f'(z_0) = \lim \frac{f(z) - f(z_0)}{z - z_0s}
    \end{equation}
    $f'(z_0)$ is the complex derivative at $z=z_0$. Of course, $D$ is open in $\mathbb{C}$ and $f: D \rightarrow \mathbb{C}$.  
\end{definition}

Of course, 

\begin{enumerate}
    \item $(\alpha f + \beta g)' = \alpha f' + \beta g'$
    \item $(fg)' = f'g + fg'$
    \item Quotient rule is the same, if $g(z_0) \neq 0$.
\end{enumerate}

\section*{Cauchy-Riemann equations}

Suppose that $f'(z_0)$ exists. Then, (recall that $f = u + iv$) 

\begin{equation}
    f'(z_0) = u_x + i v_x = v_y - iu_x, \text{approaching limit from real and complex axis}
\end{equation}

Therefore, the Cauchy-Riemann equations. 

\begin{equation}
    u_x = v_y, \quad u_y = -v_x
\end{equation}



\end{document}